% Tradução para leitura — NÃO para submissão
\documentclass[12pt,a4paper]{article}
\usepackage[utf8]{inputenc}
\usepackage[T1]{fontenc}
\usepackage[brazil]{babel}
\usepackage{mathptmx}
\usepackage[margin=2.5cm]{geometry}
\usepackage{setspace}
\usepackage{booktabs}
\usepackage{amsmath}
\usepackage{hyperref}
\onehalfspacing

\hypersetup{colorlinks=true, linkcolor=blue!70!black, urlcolor=blue!70!black}

\title{\textbf{Mesmo Prompt, Resposta Diferente: Expondo a Ilusão de Reprodutibilidade em APIs de Modelos de Linguagem de Grande Escala}}

\author{Lucas Rover$^{1,*}$, Eduardo Tadeu Bacalhau$^{2}$,\\
Anibal Tavares de Azevedo$^{3}$, Yara de Souza Tadano$^{1}$\\[6pt]
\small $^{1}$Programa de Pós-Graduação em Engenharia Mecânica, UTFPR, Ponta Grossa, PR\\
\small $^{2}$Grupo de Pesquisa em Tecnologia Aplicada à Otimização (GTAO), UFPR, Curitiba, PR\\
\small $^{3}$Faculdade de Ciências Aplicadas (FCA), Unicamp, Limeira, SP\\[4pt]
\small $^{*}$Autor correspondente: lucasrover@utfpr.edu.br}

\date{}

\begin{document}
\maketitle

\begin{center}
\textit{Tradução para leitura do manuscrito submetido à Nature Machine Intelligence.\\
Versão oficial em inglês: \texttt{nature\_mi\_main.tex}}
\end{center}

\vspace{1em}

%%=============================================================%%
%% RESUMO
%%=============================================================%%
\begin{abstract}
\noindent O mesmo prompt enviado duas vezes a uma API de modelo de linguagem de grande escala (LLM) sob configurações documentadas como ``determinísticas'' pode retornar respostas diferentes --- e essa variação é invisível para os usuários.
Neste trabalho, reportamos 4.104 experimentos controlados em oito modelos e cinco provedores de API, demonstrando que, sob decodificação greedy com temperatura zero e sementes fixas, modelos servidos por API reproduzem suas próprias saídas apenas 22,1\% das vezes (taxa de correspondência exata, EMR), enquanto modelos implantados localmente alcançam 95,6\% --- uma diferença superior a quatro vezes.
O não-determinismo persiste em fluxos de trabalho multi-turno e de geração aumentada por recuperação (RAG), nos quais um modelo obtém zero correspondências exatas em 50 execuções, porém permanece oculto porque as saídas são semanticamente equivalentes (BERTScore F1~$>$~0,97).
Um experimento de quase-isolamento demonstra que a complexidade da infraestrutura de produção --- e não a implantação em nuvem em si --- é responsável pela variação.
Fornecemos um protocolo leve de proveniência ($<$1\% de overhead) que torna essa variação detectável, levantando uma preocupação de confiabilidade para o uso crescente de LLMs em medicina, ciências físicas e análise automatizada de dados.
\end{abstract}

\textbf{Palavras-chave:} Reprodutibilidade, Modelos de linguagem de grande escala, Não-determinismo, Proveniência, Inferência por API

\vspace{1.5em}

%%=============================================================%%
%% PARÁGRAFOS DE ABERTURA (sem título de seção)
%%=============================================================%%

Modelos de linguagem de grande escala (LLMs) estão rapidamente se tornando instrumentos de pesquisa padrão.
Eles codificam conhecimento clínico, auxiliam fluxos de trabalho diagnósticos, aceleram revisão de literatura e extração de dados, e estão cada vez mais integrados à revisão por pares em larga escala.
Um corpo crescente de trabalhos emprega LLMs como avaliadores de comunicação humana, como ferramentas de avaliação psicométrica e como componentes de pipelines científicos multiagente.
Quando pesquisadores utilizam modelos baseados em API, eles confiam que configurações idênticas produzirão saídas idênticas.
A documentação das APIs reforça essa confiança ao oferecer um parâmetro de temperatura configurado como zero para comportamento ``determinístico'' --- um parâmetro que, por definição, deveria colapsar a distribuição de probabilidade de saída para um único token a cada passo.
Mas isso é uma ilusão.

A crise de reprodutibilidade mais ampla na ciência é bem documentada: mais de 70\% dos pesquisadores falharam em reproduzir o experimento de outro cientista, e o problema é particularmente agudo em IA, onde apenas 6\% dos artigos fornecem informações suficientes para reprodução.
O campo enfrenta o que tem sido chamado de sua própria ``crise de reprodutibilidade'', com preocupações crescentes de que a IA possa ativamente piorar o problema em diversas disciplinas.
Somente o vazamento de dados já foi documentado em 17 subcampos de aprendizado de máquina, e periódicos líderes em IA começaram a adotar mecanismos estruturados de reprodutibilidade em resposta.
No entanto, apesar dessa atenção à reprodutibilidade em tempo de treinamento, a reprodutibilidade da \emph{inferência} --- a implantação real de modelos para gerar saídas --- permanece amplamente inexaminada.
Essa lacuna é particularmente consequente porque a inferência é o estágio em que LLMs interagem com dados de pesquisa e produzem as saídas que entram em registros científicos, documentos de políticas e fluxos de trabalho clínicos.
Ferramentas existentes de rastreamento de experimentos, como o MLflow, foram projetadas para pipelines de treinamento e métricas numéricas, não para proveniência de geração de texto em tempo de inferência, deixando pesquisadores sem ferramentas adequadas para detectar ou documentar variação em tempo de inferência.

Conduzimos 4.104 experimentos controlados em oito modelos e cinco provedores de API independentes e descobrimos que, mesmo sob temperatura~$=$~0 com sementes aleatórias fixas, modelos servidos por API reproduzem suas próprias saídas apenas 22,1\% das vezes, enquanto modelos locais de pesos abertos alcançam 95,6\% de reprodutibilidade exata --- uma diferença superior a quatro vezes (IC bootstrap 95\% para a razão: 2,48--3,61$\times$).
Esse não-determinismo é inteiramente invisível para os usuários: as saídas são semanticamente equivalentes (BERTScore F1~$>$~0,97), mas textualmente diferentes, o que significa que pesquisadores não conseguem detectar a variação simplesmente lendo as saídas.
O fenômeno não se limita a consultas únicas.
À medida que LLMs avançam para pipelines agênticos que projetam experimentos e geram hipóteses, essa variação oculta se acumula a cada passo.
Editoriais recentes sublinharam a urgência de transparência em sistemas de IA multiagente, no entanto, a reprodutibilidade de linha de base dos componentes individuais --- as próprias chamadas de API --- nunca foi sistematicamente quantificada.

Aqui, fazemos quatro contribuições.
Primeiro, revelamos e quantificamos o não-determinismo oculto em todos os cinco principais provedores de API testados (OpenAI, Anthropic, Google, DeepSeek, Perplexity), demonstrando que é uma propriedade geral de APIs de LLM servidas em nuvem, e não uma anomalia específica de um provedor.
Segundo, mostramos que esse não-determinismo persiste --- e se amplifica --- em fluxos de trabalho multi-turno e de geração aumentada por recuperação (RAG), nos quais o Claude Sonnet~4.5 obtém zero correspondências exatas em 50 execuções de RAG.
Terceiro, demonstramos que a implantação em nuvem \emph{per se} não causa não-determinismo: a mesma arquitetura LLaMA~3 8B servida pelo endpoint em nuvem da Together AI alcança taxas de correspondência exata próximas às locais, implicando a complexidade da infraestrutura de produção (paralelismo de tensores, decodificação especulativa, batching dinâmico) em vez do meio em nuvem.
Quarto, fornecemos um protocolo leve e de padrão aberto de proveniência --- fundamentado no W3C PROV e estendendo o paradigma de documentação estruturada de Model Cards e Datasheets --- que adiciona menos de 1\% de overhead e torna essa variação invisível visível, auditável e atribuível.


%%=============================================================%%
%% RESULTADOS
%%=============================================================%%
\section*{Resultados}

\subsection*{Modelos baseados em API falham em reproduzir saídas sob configurações determinísticas}

Avaliamos oito modelos em duas tarefas centrais --- extração estruturada (saída JSON) e sumarização científica (texto livre) --- sob decodificação greedy (temperatura~$=$~0) com sementes aleatórias fixas, a configuração que a documentação das APIs apresenta como determinística.
Os resultados revelam uma divisão marcante (Tabela~1, Fig.~1).

Modelos implantados localmente alcançam reprodutibilidade bitwise quase perfeita a perfeita.
Gemma~2 9B produz EMR~$=$~1,000 [1,00;\,1,00] em todas as tarefas e condições --- cada saída é idêntica caractere por caractere em cinco repetições, confirmado por comparação de hash SHA-256.
LLaMA~3 8B atinge EMR~$=$~0,987 [0,96;\,1,00] para extração estruturada e 0,947 [0,89;\,0,99] para sumarização; Mistral 7B alcança 0,960 [0,88;\,1,00] e 0,840 [0,72;\,0,96], respectivamente.
Os desvios menores em LLaMA~3 e Mistral são inteiramente atribuíveis a um efeito de cold-start: análise post-hoc revela que, em todos os sete grupos de abstracts não unânimes, a primeira repetição era o único outlier --- provavelmente devido à inicialização do estado de cache da GPU no primeiro forward pass.
Descartar uma única inferência de aquecimento resulta em EMR~$=$~1,000 para ambos os modelos (Tabela de Dados Estendidos~5).

Modelos servidos por API contam uma história drasticamente diferente.
Sob a mesma decodificação greedy com sementes fixas, GPT-4 alcança EMR~$=$~0,443 [0,32;\,0,57] para extração e 0,230 [0,16;\,0,30] para sumarização.
Claude Sonnet~4.5 alcança 0,190 [0,05;\,0,40] e 0,020 [0,00;\,0,05], respectivamente --- o último significando que, em 50 comparações pareadas, efetivamente nenhum par de saídas de sumarização era caractere-idêntico.
Perplexity Sonar cai ainda mais para EMR~$=$~0,100 para extração e 0,010 para sumarização.
A média geral local de EMR é 0,956 versus 0,221 para modelos de API --- significando que, em média, menos de um em cada quatro pares de saídas de API são idênticos sob condições documentadas como determinísticas.
Essa diferença sobrevive à correção de Holm--Bonferroni em 68 testes de hipótese (51 significativos em $\alpha$~$=$~0,05; Seção Suplementar~S9) e é confirmada por grandes tamanhos de efeito de Cliff's delta ($\delta$~$=$~0,784--0,896), teste $t$ pareado (Cohen's $d$~$>$~1,6) e uma análise de subamostra balanceada de 10 abstracts que controla diferenças de tamanho amostral entre modelos (EMR local~$=$~0,953 vs.\ EMR API~$=$~0,304, diferença de 3,1$\times$; Tabela de Dados Estendidos~4).
Um experimento controle de formato de chat (200 execuções) confirmou que o formato do prompt (completion versus chat API) não contribui para a variação observada (Seção Suplementar~S7).


\subsection*{O não-determinismo varia substancialmente entre provedores}

Os cinco provedores de API abrangem uma ampla faixa de reprodutibilidade, com uma diferença de 80 vezes entre o mais e o menos reprodutível (Fig.~1).
DeepSeek Chat alcança a maior reprodutibilidade de API (EMR~$=$~0,800 para extração, 0,760 para sumarização), aproximando-se do desempenho de modelos locais.
GPT-4 ocupa uma posição intermediária (EMR~$=$~0,443 para extração, 0,230 para sumarização), com a diferença extração--sumarização sugerindo que a estrutura de saída da tarefa media a reprodutibilidade --- saídas restritas a JSON deixam menos espaço para variação lexical do que resumos em texto livre.
Claude Sonnet~4.5 (0,190/0,020) e Perplexity Sonar (0,100/0,010) ocupam a extremidade inferior, enquanto Gemini~2.5 Pro (avaliado apenas em tarefas multi-turno e RAG) alcança EMR~$=$~0,010--0,070.

Essa variação entre APIs é notável porque todos os provedores expõem os mesmos parâmetros ``determinísticos'' voltados ao usuário (temperatura zero, semente fixa quando suportado).
A variação, portanto, reflete diferenças na infraestrutura de serviço de produção que são invisíveis aos usuários --- consistente com chamadas editoriais para rastreamento explícito de versão de modelo e divulgação de prompt em frameworks de pesquisa baseados em LLM.
A diferença extração--sumarização é em si informativa: a extração JSON estruturada restringe o espaço de saída mais fortemente do que a sumarização em texto livre, deixando menos sequências de tokens viáveis e, portanto, menos oportunidade para divergência não-determinística.
Essa dependência de tarefa sugere que o impacto prático do não-determinismo variará entre aplicações de pesquisa, com tarefas de geração aberta (escrita criativa, formulação de hipóteses) provavelmente exibindo variação ainda maior.

A reprodutibilidade particularmente baixa do Perplexity Sonar reflete sua arquitetura aumentada por busca, onde a recuperação web em tempo real introduz uma fonte adicional de variação que se acumula com o não-determinismo interno do modelo.
Essa observação é particularmente relevante à medida que abordagens aumentadas por recuperação se tornam padrão em aplicações científicas: a combinação de geração não-determinística com recuperação variável cria um desafio multiplicativo de reprodutibilidade que pesquisadores devem considerar.


\subsection*{A implantação em nuvem não impede a reprodutibilidade}

Uma hipótese natural é que a implantação em nuvem em si --- latência de rede, balanceamento de carga, hardware compartilhado --- causa o não-determinismo.
Para testar isso, projetamos uma sonda de quase-isolamento: avaliamos a mesma arquitetura LLaMA~3 8B servida pelo endpoint em nuvem da Together AI (quantização INT4) sob prompts, sementes e temperatura idênticos à implantação local.
Se a hospedagem em nuvem fosse o fator determinante, essa implantação deveria exibir não-determinismo similar ao de APIs.

Isso não ocorre.
O LLaMA~3 servido em nuvem alcança EMR~$=$~1,000 [1,00;\,1,00] para extração e 0,880 [0,70;\,1,00] para sumarização --- praticamente idêntico à implantação local no mesmo subconjunto de 10 abstracts (EMR~$=$~1,000 e 0,920, respectivamente; intervalos de confiança bootstrap 95\% se sobrepõem).
Esse resultado fornece evidências de que a implantação em nuvem \emph{per se} não causa não-determinismo; ao contrário, a variabilidade substancial observada em GPT-4, Claude e Gemini é consistente com a complexidade da infraestrutura de serviço de produção.

Seis mecanismos bem documentados em inferência distribuída em GPU podem independentemente produzir saídas não-determinísticas mesmo sob decodificação greedy (ver Métodos): aritmética de ponto flutuante não-associativa, acumulação de precisão mista em BF16/FP16, paralelismo de tensores multi-GPU, não-determinismo do kernel FlashAttention, batching dinâmico com escalonamento contínuo de requisições e decodificação especulativa.
Nossa implantação local em GPU única elimina os mecanismos~3--6 inteiramente, e a aritmética inteira GGML Q4 mitiga o mecanismo~2 --- explicando a reprodutibilidade local quase perfeita como uma consequência prevista do ambiente de execução mais simples.
Importantemente, o resultado da Together AI mostra que esta não é uma limitação inerente da implantação em nuvem: um ambiente de serviço em nuvem bem controlado pode alcançar reprodutibilidade próxima à local.
A implicação para provedores é que modos de execução determinísticos são tecnicamente viáveis, embora possam envolver trade-offs de desempenho que as arquiteturas de serviço atuais priorizam de forma diferente.


\subsection*{Fluxos de trabalho complexos amplificam a lacuna de reprodutibilidade}

Aplicações modernas de LLM raramente consistem em chamadas de API únicas.
Diálogos de refinamento multi-turno e pipelines de geração aumentada por recuperação (RAG) são cada vez mais comuns em fluxos de trabalho de pesquisa.
Para avaliar se o não-determinismo se acumula nessas configurações, avaliamos cinco modelos em uma tarefa de refinamento de três turnos (extrair, receber feedback, refinar) e uma tarefa de extração RAG (extrair campos estruturados de um abstract com uma passagem de contexto recuperado preposta).

Modelos locais mantêm alta reprodutibilidade mesmo sob esses regimes mais complexos (Fig.~2).
Gemma~2 9B e Mistral 7B alcançam EMR~$=$~1,000 perfeito para multi-turno e RAG.
LLaMA~3 8B mostra EMR~$=$~0,880 [0,76;\,1,00] para multi-turno e 0,960 [0,88;\,1,00] para RAG --- ligeiramente inferior ao turno único, consistente com propagação de erro entre turnos de diálogo onde cada resposta condiciona a próxima.

Ambos os modelos de API exibem reprodutibilidade próxima de zero.
Claude Sonnet~4.5 alcança EMR~$=$~0,040 [0,00;\,0,08] para multi-turno e 0,000 [0,00;\,0,00] para RAG --- em 50 execuções, nenhum par de saídas de RAG era caractere-idêntico, com uma distância de edição normalizada (NED) média de 0,256 indicando variação substancial no nível de superfície.
Gemini~2.5 Pro, apesar de suportar um parâmetro de semente, alcança EMR~$=$~0,010 [0,00;\,0,03] para multi-turno e 0,070 [0,02;\,0,13] para RAG.
A convergência de dois provedores de API independentes --- usando diferentes arquiteturas de modelo, pilhas de serviço e infraestrutura geográfica --- em reprodutibilidade multi-turno próxima de zero fornece forte evidência de que esta é uma propriedade geral da inferência servida em nuvem para regimes de interação complexa, não um artefato de qualquer provedor único.
Essa descoberta é particularmente preocupante para o paradigma emergente de fluxos de trabalho de IA agêntica, onde LLMs são encadeados em múltiplos passos em pipelines automatizados.
Se cada passo introduz variação não-determinística independente, o efeito composto ao longo de um pipeline de múltiplos passos tornaria a reprodutibilidade de ponta a ponta essencialmente inatingível sem rastreamento explícito de proveniência em cada nó.


\subsection*{As saídas divergem textualmente mas não semanticamente}

Se as saídas de API variassem aleatoriamente --- produzindo resultados sem sentido ou contraditórios entre execuções --- o problema seria fácil de detectar.
Em vez disso, nossa análise métrica multinível revela um padrão mais insidioso (Fig.~3).
Em todos os modelos e condições, o BERTScore F1 permanece acima de 0,97 mesmo quando o EMR se aproxima de zero.
Os scores ROUGE-L permanecem igualmente altos ($>$0,85 para todos os modelos), confirmando sobreposição substancial no nível de tokens.
Essa dissociação em três níveis --- baixa identidade bitwise (EMR), divergência superficial moderada (NED~$=$~0,07--0,24), alta preservação semântica (BERTScore~$>$~0,97) --- define o fenômeno como não-determinismo ``oculto'': mesmo significado, palavras diferentes.
Um pesquisador lendo duas saídas do mesmo prompt as julgaria equivalentes; apenas comparação sistemática revela que são textualmente distintas.

No entanto, essa variação oculta tem consequências práticas que se estendem além das diferenças superficiais.
Uma análise de divergência no nível de campos das saídas de extração estruturada do GPT-4 revela que 100\% dos pares de saídas não-idênticos (24 de 30 grupos de abstracts) diferem em pelo menos um campo relevante para conclusões --- objetivo, método ou resultado-chave (Tabela de Dados Estendidos~7).
O campo \texttt{key\_result} diverge em 67\% dos grupos, e \texttt{method} em 57\%, significando que a variação textual não se limita a formatação ou palavras de preenchimento, mas afeta o conteúdo substantivo que pesquisadores usariam para análises downstream.
Para pipelines automatizados de síntese de evidências, onde as saídas são analisadas programaticamente em vez de lidas por humanos, mesmo diferenças lexicais menores podem se propagar para diferentes valores extraídos, diferentes estatísticas agregadas e, ultimamente, diferentes conclusões.

As implicações se estendem além de estudos individuais.
Trabalhos recentes mostraram que LLMs usados como avaliadores de comunicação empática são sensíveis a variações sutis de entrada, e avaliações psicométricas do comportamento de LLM revelam padrões de consistência que dependem da arquitetura do modelo e da abordagem de fine-tuning.
Essas descobertas ressaltam que o não-determinismo não é meramente uma curiosidade técnica, mas uma propriedade que interage com a forma como LLMs são implantados como instrumentos de medição científica --- e uma que os frameworks de avaliação atuais não consideram.
Para o corpo crescente de trabalhos que usa LLMs como anotadores, avaliadores ou extratores de dados no lugar de participantes humanos, nossos resultados indicam que o ``instrumento'' em si introduz ruído de medição que não é caracterizado nem reportado nas seções de metodologia padrão.


\subsection*{Paradoxo da temperatura: decodificação greedy não é greedy para APIs}

A temperatura é o principal parâmetro controlável pelo usuário que governa a variabilidade de saída.
Para modelos locais, ela se comporta como a teoria prevê: sob a varredura de temperatura (Fig.~4), modelos locais mostram um declínio monotônico limpo de EMR quase perfeito em $t$~$=$~0 para EMR~$\approx$~0 em $t$~$=$~0,7.
Isso confirma que a decodificação greedy local é verdadeiramente greedy --- colapsando a distribuição de saída para uma única sequência determinística.
Modelos de API, no entanto, partem de uma linha de base já baixa em $t$~$=$~0 e mostram um padrão mais complexo que viola a monotonicidade esperada.

Claude Sonnet~4.5 exibe uma resposta não-monotônica particularmente anômala: o EMR de extração \emph{aumenta} de 0,067 em $t$~$=$~0,0 para 0,700 em $t$~$=$~0,3 antes de declinar para 0,133 em $t$~$=$~0,7.
Esse resultado contraintuitivo --- onde adicionar aleatoriedade \emph{melhora} a reprodutibilidade --- sugere que o caminho de decodificação $t$~$=$~0 da Anthropic pode expor mais estocasticidade no nível de infraestrutura (variação de ponto flutuante de hardware, roteamento de requisições) do que uma temperatura positiva pequena que ativa um caminho de amostragem mais estável.
GPT-4 mostra um padrão menos dramático mas qualitativamente similar: a diferença entre $t$~$=$~0 e $t$~$=$~0,3 é menor para modelos de API do que para modelos locais, porque a linha de base da API já está longe de ser determinística.
Em $t$~$=$~0,7, modelos locais e de API convergem para EMR uniformemente baixo, mas por trajetórias diferentes: modelos locais declinam monotonicamente de quase 1,0, enquanto modelos de API seguem uma curva mais plana de uma linha de base já degradada.

Esse paradoxo da temperatura tem implicações práticas importantes.
Pesquisadores que configuram a temperatura para zero esperando saídas determinísticas não estão, para modelos de API, alcançando a decodificação greedy que pretendem.
A relação temperatura--reprodutibilidade para modelos de API depende de detalhes de implementação específicos do provedor que são opacos aos usuários e não documentados nas referências das APIs, tornando impossível para pesquisadores prever ou controlar o grau de não-determinismo em seus experimentos sem medição empírica.


%%=============================================================%%
%% DISCUSSÃO
%%=============================================================%%
\section*{Discussão}

Nossos achados revelam uma lacuna de reprodutibilidade pervasiva e previamente invisível na pesquisa baseada em LLM.
Sob as configurações que a documentação das APIs apresenta como determinísticas --- temperatura zero, semente fixa --- modelos servidos por API falham em reproduzir suas próprias saídas aproximadamente quatro em cada cinco vezes.
Essa lacuna persiste em cinco provedores independentes, se estende a fluxos de trabalho multi-turno e RAG, e sobrevive a correção estatística rigorosa.
A implicação imediata é que uma porção substancial da pesquisa publicada que depende de LLMs baseados em API pode conter resultados não-reprodutíveis sem o conhecimento dos autores, uma preocupação amplificada pelo uso crescente de LLMs em revisão por pares, geração de hipóteses e pipelines agênticos onde o não-determinismo se acumula a cada passo.

Trabalhos anteriores documentaram o não-determinismo de LLM em contextos isolados --- saídas inconsistentes em benchmarks de PLN, falha das configurações ``determinísticas'' de API e geração de código não-determinística.
Nosso estudo avança além dessas observações ao comparar sistematicamente implantações locais versus de API em cinco provedores, estender as medições a fluxos de trabalho multi-turno e RAG, e fornecer um framework de atribuição causal que distingue complexidade de infraestrutura de implantação em nuvem.

O resultado de quase-isolamento da Together AI fornece uma percepção mecanística chave: como a mesma arquitetura de modelo alcança reprodutibilidade próxima à local via uma API em nuvem diferente, o não-determinismo observado em GPT-4, Claude e Gemini é atribuível à complexidade do serviço de produção --- não à implantação em nuvem em si.
Isso sugere que a reprodutibilidade é uma escolha de design que provedores poderiam endereçar através de modos de execução determinísticos ou documentação transparente, em vez de uma limitação inerente.
O estado atual --- onde temperatura zero é documentada como ``determinística'' sem divulgar variação no nível de infraestrutura --- cria uma falsa sensação de reprodutibilidade que pode ser mais danosa do que estocasticidade abertamente reconhecida.

Essas descobertas também têm implicações regulatórias.
O EU AI Act requer rastreabilidade para sistemas de IA de alto risco, e o NIST AI RMF enfatiza transparência; o não-determinismo oculto compromete diretamente ambos.
Diretrizes éticas para uso de LLM em escrita acadêmica similarmente assumem que configurações documentadas produzem saídas previsíveis --- uma suposição que nossos resultados mostram ser injustificada.

Embora a preservação semântica que observamos (BERTScore F1~$>$~0,97) sugira que a qualidade da saída não é corrompida, nossa análise no nível de campos mostra que mesmo saídas ``semanticamente equivalentes'' diferem em campos relevantes para conclusões; a taxa de 78\% de não-correspondência sob configurações ``determinísticas'' representa um obstáculo sério para submissões regulatórias, revisões sistemáticas e pipelines automatizados.

O protocolo de proveniência que fornecemos oferece um caminho prático adiante.
Fundamentado no W3C PROV e adicionando menos de 1\% de overhead (Tabela de Dados Estendidos~6), ele cria registros auditáveis vinculando cada saída ao seu contexto de geração.
Sua propriedade chave é o \emph{diagnóstico diferencial}: quando todos os hashes coincidem exceto o hash de saída, a divergência é automaticamente atribuída ao processo de geração --- um resultado confirmado em todas as 4.104 execuções (Seção Suplementar~S4).
Isso implementa o que editoriais recentes têm advogado e estende Model Cards e Datasheets para a camada de inferência.
Propomos um padrão mínimo de reporte --- identidade do modelo, parâmetros, métricas de reprodutibilidade, modo de implantação e hashes de saída --- que nosso protocolo automatiza a custo negligível (Informação Suplementar).

Nosso estudo possui diversas limitações.
Cobre oito modelos em quatro tarefas, mas exclui geração de código, raciocínio matemático e escrita criativa.
Os dados de entrada compreendem 30 abstracts de IA/ML em inglês; outros idiomas e domínios podem mostrar efeitos amplificados.
Os experimentos com GPT-4 usaram o snapshot \texttt{gpt-4-0613}; versões mais novas podem diferir, embora os mecanismos no nível de infraestrutura que identificamos sejam gerais.
Os experimentos multi-turno e RAG incluem apenas dois provedores de API; estender a outros fortaleceria a generalidade.
Os tamanhos amostrais (10--30 abstracts por modelo) são bem dimensionados para os grandes efeitos observados ($d$~$>$~1,6), mas podem perder fenômenos mais sutis.


%%=============================================================%%
%% CONCLUSÕES
%%=============================================================%%
\section*{Conclusões}

Demonstramos que as configurações ``determinísticas'' oferecidas pelos principais provedores de API de LLM não entregam determinismo na prática, criando uma lacuna de reprodutibilidade oculta que afeta qualquer estudo que dependa de modelos servidos por API.
A lacuna é pervasiva --- abrangendo cinco provedores independentes --- persistente em fluxos de trabalho multi-turno e de geração aumentada por recuperação, e invisível aos usuários porque a equivalência semântica mascara a divergência textual.
Ao executar o mesmo modelo de pesos abertos localmente e via API em nuvem, isolamos a causa: a complexidade da infraestrutura de produção, e não a implantação em nuvem em si, impulsiona o não-determinismo.
Nosso protocolo leve de proveniência, fundamentado no W3C PROV e requerendo menos de 1\% de overhead, oferece um mecanismo prático para tornar essa variação visível, auditável e atribuível.
Juntamente com o padrão mínimo de reporte que propomos, ele fornece um caminho acionável para pesquisa reprodutível baseada em LLM.

Essas descobertas levantam uma preocupação de \textbf{confiabilidade} mais ampla.
Modelos de linguagem de grande escala não estão mais confinados à geração de texto; eles estão cada vez mais incorporados em suporte à decisão clínica, monitoramento ambiental, projeto de engenharia, coleta e análise automatizada de dados, e outros domínios onde as saídas informam diretamente decisões com consequências no mundo físico.
Nesses contextos, a incapacidade de reproduzir ou verificar a saída de um modelo não é meramente um inconveniente acadêmico, mas uma ameaça à confiabilidade dos sistemas que dependem dele.
Nossos resultados demonstram que, sem rastreamento explícito de proveniência, a confiabilidade de qualquer pipeline baseado em LLM construído sobre inferência por API não pode ser assumida --- ela deve ser medida, documentada e continuamente monitorada.
À medida que a integração da IA na ciência e na sociedade se acelera, estabelecer a infraestrutura para saídas de modelos verificáveis e reprodutíveis não é uma sutileza técnica, mas uma necessidade científica e social.

\end{document}
